%%%PAGE 185 rot=0 legibility=fair summary=粘贴印刷纸（盖住笔记页中段）：滑块小车(2)问、传送带弹簧动量题23、气体V-T圆题24，均带手写圈画与草算；纸外上沿/下沿笔迹与p186为同一笔记页，未重复转录
% 说明：本页是 p186 同一笔记页的另一张照片，其中段被一张粘贴的印刷纸（紫色）盖住，因此：
%  (1) 纸外上沿（验证力的平行四边形定则：“不用让OA⊥OB，不用让O点不变，可以换成弹性绳 (只需稳定即可) (分力和合力可在同一个图示中表示出来) (不影响受力的方向和钩码个数)”及钩码链图）
%      与下沿（自由落体验机守：装置图，有两处不当 ①用交流电不用直流电 ②重物离打点计时器太远，v^2=2gh，(1/2)v^2=gh，图线斜率表示重力加速度g）的手写内容，
%      已在 p186.tex 完整转录，此处不重复；
%  (2) 本页独有的新内容仅为粘贴印刷纸上的三道题（22题(2)问、23题、24题(1)）及其上的手写圈画/草算，转录如下。汇编时请勿把 p185 上沿/下沿与 p186 重复收录。
%%%BLOCK id=185-01 ch=03 sec=板块模型·滑块与小车 kind=problem alt=01 cont=prev y=0.12-0.41
% 粘贴的印刷题（紫色纸），右端被页边截断；原稿对“相同的恒力F”圈出，对“v_0=4m/s”“要保证滑块不能从车的左端掉下”“恒力F大小”及整行有划线
\begin{problem}[粘贴印刷题·第(2)问]
(2)已知滑块与车面间动摩擦因数 $\mu=0.2$，滑块质量 $m=1\unit{kg}$，车长 $L=2\unit{m}$，车速 $v_0=4\unit{m/s}$，取 $g=10\unit{m/s^2}$，当滑块……% 此行右端被页边截断
车面中点的同时对该滑块施加一个与车运动方向相同的恒力 $F$，要保证滑块不能从车的左端掉下，恒力 $F$ 大小……% 此行右端被页边截断
满足什么条件？

\medskip
\noindent\textbf{手写草算（右侧页边，部分被截断）：}
$2+F=a_1$；\quad $a_1t_1=4$；\quad $4t-\dfrac12$\struck{$(2+F)$}\unclear{…}；\quad $4t-2$\unclear{…}；\quad 再下方：$\dfrac{a_1}{2}$\unclear{…}，$a_1$\unclear{…}% CHECK: 右侧被页边截断，后半不可见
\end{problem}
%%%END
%%%BLOCK id=185-02 ch=07 sec=动量与能量综合·弹簧传送带 kind=problem alt=03,06 cont=none y=0.23-0.62
% 粘贴的印刷题（紫色纸），右端被页边截断；原稿圈画“足够长的”“光滑”“v=4m/s的”“顺时针”“恰好”，划线“m_A=1kg,m_B=2kg”“运动过程”“B脱离弹”“滑上传送带”“传送带顶端”“动摩擦因数μ=0.5”
\begin{problem}[粘贴印刷题 23]
23．如图所示，足够长的光滑平台上有静止的小滑块 A、B，$m_A=1\unit{kg}$，$m_B=2\unit{kg}$，两滑块之间有一段轻质……% 此行右端被页边截断
刚好处于原长，滑块 A 与轻弹簧栓接，滑块 B 未栓接弹簧，平台右端与倾角 $\theta=37^\circ$ 的倾斜传送带平滑连接，……% 此行右端被页边截断
带以 $v=4\unit{m/s}$ 的恒定速度顺时针转动。现给滑块 A 瞬时向右的冲量 $I=9\unit{N\cdot s}$，此后运动过程中，滑块 B 脱离弹……% 此行右端被页边截断
滑上传送带，并恰好能到达传送带顶端。已知滑块 B 与传送带之间的动摩擦因数 $\mu=0.5$，取重力加速度 $g=10\unit{m/s^2}$，$\sin37^\circ=0.6$，$\cos37^\circ=0.8$。求：

（1）滑块 B 刚滑上传送带时的速度大小；

（2）滑块 B 在传送带上运动过程中，摩擦力对滑块 B 做的总功；

（3）滑块 B 返回平台与滑块 A、弹簧发生二次作用过程中，弹簧的最大弹性势能（结果可用根号表示）。

%FIG p185_a 0.04 0.505 0.37 0.598
\notefig[0.33]{figures/p185_a.png}

\medskip
\noindent\textbf{手写旁注：}（2）问末尾圈出 \fbox{$\vec{f}\,\vec{x}$}（$f$、$x$ 上带箭头），其右 $9=2V$% CHECK: 首字符可能是 a/q
；图右侧草算：8N，$-8$，4，8，$32-8$，24；（3）问右侧 $\dfrac{4.5-4}{10}$
\end{problem}
%%%END
%%%BLOCK id=185-03 ch=16 sec=理想气体·V-T图像 kind=problem alt=none cont=none y=0.60-0.745
% 粘贴的印刷题（紫色纸），右端被页边截断；只有题干，无解答
\begin{problem}[粘贴印刷题 24(1)]
24．(1)如图所示，一定质量的理想气体由状态 $a$ 开始，经历 $ab$、$bc$、$cd$、$da$ 四个过程回到原状态，其 $V-T$……% 此行右端被页边截断
刚好是一个圆，气体在 $a$、$b$、$c$、$d$ 四个状态的体积分别为 $V_a=V_c$、$V_b=V_2$、$V_d=V_1$，温度分别为 $T_a=T_1$、$T_b=$……% 此行右端被页边截断
$T_c=T_2$，压强分别为 $p_a$、$p_b$、$p_c$、$p_d$，其中 $V_2=5V_1$，$T_2=5T_1$。下列说法正确的是（\qquad）

\medskip
\noindent\textbf{手写旁注（右下，下半被粘贴纸边缘遮挡）：}\unclear{$a=3V$}；\unclear{$\dfrac{10}{\ldots}$}，\unclear{.2}% CHECK: 被遮挡，可能是 V_a=3V_1 或 9=3V（与上方 9=2V 同类）
\end{problem}
%%%END
%%%PAGEEND 185
